% methodology.tex
% Methodology section for VICTOR paper.
% Intended to be \input{} from the main paper .tex file.



\section{Methodology}
\label{sec:methodology}

\begin{figure*}[t]
    \centering
    \includegraphics[width=0.98\textwidth]{figures/blicl_method.pdf}
    \caption{\textbf{\method method overview (placeholder).} At deployment,
    the demonstration pool $\mathcal{D}_\ell$ feeds two in-context
    estimators. The IC-VFE ($\phi$) conditions on one context demonstration
    to produce a calibrated progress estimate $\hat v_t$ for the current
    observation $o_t$; $\hat v_t$, fused with visual similarity, drives
    retrieval of a context chunk $c_t$ from the same pool. The IC-VLA
    ($\theta$) conditions on $c_t$ to predict the next action chunk,
    closing the loop without any parameter updates. \gap{placeholder
    figure -- replace with final diagram once the layout in
    Section~\ref{sec:context-retrieval} is finalized.}}
    \label{fig:method}
\end{figure*}

In this section we present an overview of \method: the problem setting, how
per-task context is constructed and retrieved, the architecture and
meta-training objective of the In-Context Value Function Estimator (IC-VFE),
and the architecture and meta-training objective of the IC-VLA policy.

\subsection{Problem Formulation}
\label{sec:problem-formulation}

We consider a multi-task vision-language-action setting. A task is specified by a
language instruction $\ell \in \mathcal{L}$, and an observation
$o_t = [X_t^1, \dots, X_t^n, q_t]$ consists of $n$ camera images and proprioceptive
state $q_t$. In our real-robot deployment $n=3$ and $q_t$ is an 18-dimensional
per-arm end-effector position and orientation, the orientation given as a 6D
rotation representation rather than a unit quaternion to avoid the antipodal
(double-cover) ambiguity a quaternion state introduces; the action chunk
$a_{t:t+H}$ is correspondingly 20-dimensional per step (per-arm delta
end-effector position relative to the chunk's first frame, absolute 6D
rotation, and a gripper command). A base
VLA policy $\pi_\theta(a_{t:t+H} \mid o_t, \ell)$ maps an observation and
instruction to an action chunk $a_{t:t+H}$.
Each demonstration $\tau = (o_0, a_0, \dots, o_T)$ is annotated offline
(Section~\ref{sec:dataset}) with a scalar progress
value $v_t \in [0, 1]$, with $v_0 \approx 0$ and $v_T = 1$ on success.

\paragraph{Demonstration pool.} Each meta-training task $\ell$ has an associated
demonstration set $\mathcal{D}_\ell$: individual $(o_t, a_{t:t+H})$ pairs from
$\mathcal{D}_\ell$ supply queries for the base flow-matching loss, and the same
per-task episode set is flattened into the demonstration chunk dictionary
$\mathcal{D}_\ell^{\text{ctx}}$ that both in-context estimators condition on
(Sections~\ref{sec:ic-vfe-objective}, \ref{sec:ic-vla-objective}) and that supplies
the retrieval pool defined in Section~\ref{sec:context-retrieval}. Unlike an
idealized held-out split, our current real-robot training arms draw
$\mathcal{D}_\ell^{\text{ctx}}$ from the same episodes used for training queries
rather than a disjoint subset; enforcing a strict split is future work. At
deployment there is no train/context distinction to begin with: the user supplies
a single small demonstration set $\mathcal{D}_\ell$ ($N \le 10$) for the new task,
and this entire set plays the role $\mathcal{D}_\ell^{\text{ctx}}$ played during
meta-training.

\paragraph{Bi-level structure.} We call the overall method \method\ because two
independent in-context estimators compose at inference time, both conditioned on
the same demonstration pool but in different forms (Section~\ref{sec:context-retrieval}):
(1) the IC-VFE conditions on one full context demonstration to produce a calibrated
progress estimate $\hat{v}_t \in [0,1]$ for the current observation, and (2) that
estimate, fused with visual similarity, drives a retrieval step that selects a context
chunk $c_t$ consumed by the IC-VLA to predict the next action chunk. Both estimators
are meta-trained on oracle progress values $v_t$; at deployment, oracle values are
replaced by $\hat{v}_t$ from the trained IC-VFE, so downstream retrieval quality is
bounded by IC-VFE accuracy.

\subsection{Preliminaries}
\label{sec:preliminaries}

Both the base VLA and the IC-VLA follow the $\pi_{0.5}$ recipe of a VLM backbone
paired with a flow-matching action expert
that predicts the continuous action chunk $a_{t:t+H}$~\cite{pi05}. Given a denoising
timestep $\eta \in [0,1]$, a noised action chunk $a_{t:t+H}^\eta = \eta\, a_{t:t+H} +
(1-\eta)\,\epsilon$ with $\epsilon \sim \mathcal{N}(0, I)$, and a velocity field
$u_\theta$ predicted by the action expert conditioned on $(o_t, \ell)$ (and, for the
IC-VLA, the retrieved context $c_t$), the flow-matching loss is
\begin{equation}
    L_{\text{flow}}\big(a_{t:t+H} \mid \cdot\,; \theta\big)
    = \mathbb{E}_{\eta, \epsilon}\Big[\big\|u_\theta(a_{t:t+H}^\eta, \eta \mid \cdot) -
    (a_{t:t+H} - \epsilon)\big\|^2\Big] ,
    \label{eq:flow-loss}
\end{equation}
which we reuse, without modification, as the action-chunk loss term for both the base
VLA fine-tune and the IC-VLA (Section~\ref{sec:ic-vla-objective}).

\subsection{Context Construction and Retrieval}
\label{sec:context-retrieval}

The demonstration pool $\mathcal{D}_\ell^{\text{ctx}}$ (or $\mathcal{D}_\ell$ at
deployment; Section~\ref{sec:problem-formulation}) feeds the two estimators in two
different forms. First, one whole demonstration $\tau_{\text{ctx}} \sim
\mathcal{D}_\ell^{\text{ctx}}$ is sampled and given to the IC-VFE as the context
sequence it conditions on to produce $\hat v_t$ (Section~\ref{sec:ic-vfe-architecture}).
Second, the pool is broken apart into a frame-level retrieval buffer from which the
IC-VLA's retrieval step selects a context chunk $c_t$. This mirrors the ``priming
vs.\ retrieval buffer'' split in RICL~\cite{ricl}, except here the same held-out set
serves both roles for both estimators, keeping the two meta-training procedures
aligned on identical context data.

\paragraph{Chunk dictionary construction.} We flatten the $N$-demonstration pool into
a \emph{demonstration chunk dictionary}: a key--value store whose keys are
DINO-v2~\cite{dinov2} embeddings of a chunk's first frame, and whose values are the
chunk contents (per-frame image, proprioception, and action). To
build the dictionary, each of the pool's demonstrations is split into overlapping
chunks of fixed length $L$ (in frames, $L=10$ in our implementation) with a
50\%-overlap stride $S = L/2$ between consecutive chunks; a demonstration of $T$
frames contributes $\lfloor (T-L)/S \rfloor + 1$ chunks, and each chunk carries its
oracle value (meta-training) or IC-VFE-estimated value (deployment). Because
pool chunks overlap 50\% within an episode, temporally-adjacent chunks often have
near-identical DINO-v2 embeddings; to prevent unconstrained top-$k$ retrieval from
collapsing onto several near-duplicate chunks from the same episode, retrieval
enforces an episode-diversity constraint, selecting at most one chunk per source
episode among the top-$k$ neighbors.

\begin{figure}[t]
    \centering
    \includegraphics[width=\linewidth]{figures/demo_dictionary.pdf}
    \caption{\textbf{Demonstration chunk dictionary construction (placeholder).}
    Illustration of splitting a demonstration into overlapping chunks of size $L$ and
    indexing each by the DINO embedding of its first frame. \gap{placeholder figure --
    replace with final diagram.}}
    \label{fig:chunk-dict}
\end{figure}

\paragraph{Retrieval.} A retrieval function
\begin{equation}
    f_{\text{retrieve}}\big(o_t, \hat{v}_t, \mathcal{P}_\ell\big) \;\rightarrow\;
    c_t = \big(n^{(1)}, \dots, n^{(k)}\big)
    \label{eq:retrieval-fn}
\end{equation}
queries the chunk dictionary $\mathcal{P}_\ell$ and selects the $k$ nearest neighbor
chunks $n^{(i)} = (o', a', v')$ (subject to the episode-diversity constraint above),
ranked by a fused score over visual similarity $d_{\text{vis}}(o_t, o')$ (DINO-v2
embedding $\ell_2$ distance, as in RICL) and value alignment $|\hat{v}_t - v'|$. Our
vision+value \method variant combines the two by independently $z$-scoring
$d_{\text{vis}}$ and $|\hat{v}_t - v'|$ across the current candidate pool -- so the
two signals, which live on very different scales, contribute comparably -- and
summing the result, $g(d_{\text{vis}}, |\hat v_t - v'|) = z(d_{\text{vis}}) +
z(|\hat v_t - v'|)$, rather than a fixed weighting such as RICL's distance-decay
$e^{-\lambda d}$. This fixed fusion is a simpler alternative to a jointly-trained
scalar fusion $g_\psi$ -- a small 2-layer MLP over $[d_{\text{vis}}, |\hat v_t -
v'|]$, trained end-to-end with the IC-VLA via a soft, temperature-annealed
relaxation of top-$k$ selection since hard top-$k$ is not itself differentiable --
which we leave to future work (Section~\ref{sec:limitations}).

\paragraph{Context window token layout.} The $k$ retrieved neighbor chunks $c_t$ are
placed in the IC-VLA's context ordered farthest-to-nearest, so that the nearest
neighbor immediately precedes the query observation tokens. Each neighbor chunk
contributes one image per input camera (all taken from the chunk's first frame)
together with a single combined text/action token span, produced by the base
$\pi_{0.5}$ tokenizer applied to the neighbor's own task string, its initial
proprioceptive state (digitized directly into the prompt text -- the same mechanism
$\pi_{0.5}$ already uses to embed the query's own state), and its
action-horizon-length action window (FAST-tokenized); this reuses $\pi_{0.5}$'s
existing text/action embedding path exactly, introducing no new token type. Each
neighbor block receives a learned neighbor-rank embedding (added to every token in
the block) so the IC-VLA can in principle distinguish nearer from farther context;
all real-robot runs reported in this work use $k=1$ (Section~\ref{sec:limitations}),
so this embedding currently only marks the presence of retrieved context rather than
distinguishing among multiple ranks. This layout gives the context-conditioned
policy
\begin{equation}
    \pi_\theta\big(a_{t:t+H} \mid o_t, \ell, c_t\big) .
    \label{eq:ic-vla-policy}
\end{equation}

\subsection{IC-VFE Architecture and Demonstration Bin Tokenization}
\label{sec:ic-vfe-architecture}

Rather than estimating value from the query observation alone (as in RECAP's
$p_\phi(V \mid o_t, \ell)$~\cite{pi06}), the IC-VFE additionally conditions on one full
context demonstration's downsampled frame sequence $\tau_{\text{ctx}}$
(Section~\ref{sec:context-retrieval}):
\begin{equation}
    p_\phi\big(V \mid o_t, \ell, \tau_{\text{ctx}}\big) \in \Delta^B ,
    \label{eq:ic-vfe}
\end{equation}
a distribution over $B = 201$ discretized value bins (following RECAP's distributional
value head), from which we recover a scalar estimate $\hat{v}_t = \sum_b p_\phi(V = b
\mid \cdot)\, \nu(b)$. Intuitively, the context demonstration lets the IC-VFE
calibrate what progress looks like for this specific task, object, or embodiment
instance before scoring the live observation, rather than relying purely on
generalization from pretraining. A no-context ablation removes only the demonstration
sequence while retaining the same current-observation encoder and prediction head.

\paragraph{Architecture.} The context demonstration is downsampled end-anchored (most
recent frames preferred, at most 200 slots) into a fixed number of temporal
\emph{context slots}, each visually tokenized by a frozen Gemma 3 vision tower and a
trainable patch merger (compressing 256 pooled visual tokens per camera to 16, and
projecting into the 640-dimensional Gemma 270M embedding space), alongside a
proprioceptive token; with 3 cameras this gives 49 tokens per slot and up to 9{,}800
context tokens total. The current observation is separately encoded at two
resolutions: a trainable detail connector projects all 256 pooled tokens per camera
independently (preserving fine visual detail), while the same patch merger used for
the context path also produces a 16-token summary (placing the current frame in the
same compressed visual space as the demonstration). Context tokens, the dual-resolution
current-observation tokens, current proprioceptive state, the tokenized instruction
(up to 32 tokens), and a learnable value-query token are concatenated and processed by
a Gemma 3 270M causal decoder; because the value-query token is last, its final hidden
state attends to all preceding tokens and is mapped by a value head to a categorical
distribution over the 201 value bins.

\paragraph{Demonstration-derived value-bin tokenization.} The value target is defined
by the relative progress of the current query frame within its successful
demonstration. For a successful trajectory $\tau = (o_0, \dots, o_{T-1})$, the
continuous value assigned to frame $t$ is
\begin{equation}
    y_t = \frac{t}{T-1}, \qquad T > 1 ,
\end{equation}
so the initial frame has value $0$, the terminal successful frame has value $1$, and
intermediate frames are labeled by linear relative progress; for the degenerate case
$T=1$, the single terminal frame is assigned $y_0 = 1$. We discretize this continuous
target into $B = 201$ uniformly spaced value bins. The scalar represented by bin $b
\in \{0, \dots, B-1\}$ is $\nu_b = b / (B-1)$, so adjacent bins are separated by
$1/200 = 0.005$. The target bin index for value $y_t$ is
\begin{equation}
    b_t = \operatorname{clip}\left(
        \operatorname{round}\left[y_t (B-1)\right],\, 0,\, B-1
    \right) .
\end{equation}
We refer to this operation as demonstration-derived value-bin tokenization because the
categorical target is obtained from the query frame's position in a successful
demonstration; these 201 bins are supervision classes for the value head, not tokens
appended to the multimodal input sequence.

\subsection{IC-VFE Meta-Learning Training Objective}
\label{sec:ic-vfe-objective}

Let $z_\theta(o_t, \ell, \tau_{\text{ctx}}) \in \mathbb{R}^{B}$ denote the logits
produced by the value head, where $B = 201$. The predicted categorical value
distribution is
\begin{equation}
    p_\theta(b \mid o_t, \ell, \tau_{\text{ctx}})
    = \frac{\exp z_{\theta, b}}{\sum_{b'=0}^{B-1} \exp z_{\theta, b'}} .
\end{equation}
IC-VFE is trained with categorical cross-entropy against the demonstration-derived
target bin $b_t$,
\begin{equation}
    \mathcal{L}_{\mathrm{IC\text{-}VFE}}(\theta)
    = -\mathbb{E}_{(o_t, \ell, \tau_{\text{ctx}}) \sim \mathcal{B}}
    \left[\log p_\theta(b_t \mid o_t, \ell, \tau_{\text{ctx}})\right] ,
\end{equation}
where $\mathcal{B}$ denotes the training distribution over query frames, instructions,
and same-task context demonstrations sampled from $\mathcal{D}_\ell^{\text{ctx}}$
(Section~\ref{sec:problem-formulation}). For the no-context ablation, the objective is
unchanged and the conditioning variable $\tau_{\text{ctx}}$ is omitted. The frozen
vision tower and fixed average-pooling operator are excluded from optimization;
gradients update the detail connector, shared patch merger, proprio encoder, learned
source embeddings, Gemma 270M backbone, value-query token, and value head.

At inference time, a scalar value is recovered as the expectation of the predicted
categorical distribution,
\begin{equation}
    \hat{y}_t = \mathbb{E}_{p_\theta}[\nu_b] = \sum_{b=0}^{B-1} p_\theta(b \mid o_t, \ell, \tau_{\text{ctx}})\, \nu_b .
\end{equation}
This distributional formulation preserves uncertainty over progress while providing a
scalar estimate in $[0, 1]$ when required. Mean-squared error and mean-absolute error
between $\hat{y}_t$ and $y_t$ may be reported as diagnostics, but they are not
auxiliary optimization terms: the training loss is the categorical cross-entropy
alone.

\subsection{IC-VLA Architecture and Context Chunk Tokenization}
\label{sec:ic-vla-architecture}

The IC-VLA shares its base architecture and fine-tuning regime with $\pi_{0.5}$
(Table~\ref{tab:ft-regime}), extended to accept the retrieved context chunk $c_t$
alongside the query observation, following the general recipe of placing retrieved
neighbors in context ordered by rank (as in RICL~\cite{ricl}). The
token-level layout of this context window --- how each retrieved neighbor is
tokenized, how many neighbors are retrieved, and how neighbor tokens are ordered
relative to the query --- is given in Section~\ref{sec:context-retrieval}.

\subsection{IC-VLA Meta-Learning Training Objective}
\label{sec:ic-vla-objective}

The IC-VLA is trained with the context-conditioned policy $\pi_\theta(a_{t:t+H}
\mid o_t, \ell, c_t)$ (Equation~\ref{eq:ic-vla-policy}): the action expert predicts
the continuous action chunk $a_{t:t+H}$ conditioned on the query observation,
instruction, and retrieved context, trained with the flow-matching loss
$L_{\text{flow}}$ (Equation~\ref{eq:flow-loss}), following the fine-tuning regime in
Table~\ref{tab:ft-regime} (Appendix~\ref{app:dataset-prep}) and mirroring the
stop-gradient interface between the VLM backbone and the action expert used for the
base $\pi_{0.5}$ fine-tune (Table~\ref{tab:ft-regime}). The training objective is
\begin{equation}
    \mathcal{L}_{\mathrm{IC\text{-}VLA}}(\theta)
    = \mathbb{E}_{(o_t, \ell, c_t) \sim \mathcal{B}}
    \Big[L_{\text{flow}}\big(a_{t:t+H} \mid o_t, \ell, c_t; \theta\big)\Big] ,
    \label{eq:ic-vla-loss}
\end{equation}
During
meta-training, the context chunk $c_t$ is sampled by $f_{\text{retrieve}}$
(Equation~\ref{eq:retrieval-fn}) using oracle progress values $v_t$ from
$\mathcal{D}_\ell^{\text{ctx}}$ (Section~\ref{sec:problem-formulation}); at
deployment, oracle values are replaced by the trained IC-VFE's estimate $\hat{v}_t$,
so retrieval quality at inference time is bounded by IC-VFE accuracy. \gap{confirm
accuracy.}

